\documentclass[12pt,a4paper,oneside]{article}
\usepackage[brazil]{babel}
\usepackage[utf8]{inputenc}
\usepackage{olymp}


\newlength{\exmpwidouf}
\exmpwidinf=10cm
\exmpwidouf=6cm


\setcounter{problem}{3}

\begin{document}

\begin{problem}{Recomendação Netflix}{netflixvetor}{1}

O programa de recomendação de filmes, séries e documentários da Netflix compara as preferências de diversos usuários para decidir quais títulos recomendar a um usuário. Neste problema, você deve escrever um programa para comparar as preferências de um usuário $U$ com as preferências de vários outros usuários e informar qual usuário mais se assemelha a $U$.

Consideramos $N$ usuários e $T$ títulos. Cada usuário indica a preferência por cada um dos títulos. As preferências são valores inteiros de 0 e 9 que representam a nota do usuário para um título. A similaridade aqui é simplificada e deve ser calculada como o número de preferências (notas) iguais entre $U$ e um outro usuário. Por exemplo, para $T = 2$ títulos, se um usuário fornece as notas 4 9 para os dois títulos e $U$ fornece as notas 4 8, a similaridade entre os dois é igual a 1. Se $U$ fornece as notas 3 8, a similaridade é igual a 0.

% \Notes

% Você pode criar uma função para o cálculo do fatorial.

\InputFile

A primeira linha contém dois inteiros, $N$ e $T$, respectivamente o número de usuários e o número de títulos (portanto, de notas). Em seguida, há uma linha com as $T$ preferências do usuário $U$. Por fim, há $N$ linhas com $T$ inteiros cada, informando as preferências de cada usuário por cada título. Restrição: $1 \leq N, T \leq 100$.

\OutputFile

Seu programa gera apenas uma linha de saída, com um valor entre $1$ e $N$, indicando o número do usuário que mais se assemelha a $U$. Caso haja empate entre mais de um usuário que mais se assemelha a $U$, deve ser escrito o menor número entre os empatados. Por exemplo, se os usuários que mais se assemelham a $U$ são os usuários 20 e 23, deve ser escrito 20.

% \Notes
% Preste atenção aos tipos das variáveis, pois $F(90) \approx 2^{61}$.

\Example

\begin{example}
\exmp{
3 5
3 7 3 3 0
3 7 4 4 4
3 4 5 6 7
3 4 4 3 0
}{
3
}%
\end{example}

\begin{example}
\exmp{
4 3
7 7 7
1 2 7
7 7 9
8 7 7
6 6 6
}{
2
}%
\end{example}

\begin{example}
\exmp{
3 6
3 8 4 4 0 9
9 3 9 7 4 9
9 6 6 2 4 3
3 0 7 6 2 3
}{
1
}%
\end{example}

\end{problem}

\end{document}
